\documentclass{article}
\usepackage{amsmath, amssymb}
\setlength{\parindent}{0pt}
\renewcommand{\arraystretch}{1.35}

\begin{document}

\textbf{CS 2050 --- Notes 09/09}

\bigskip

\textbf{Definitions}

\medskip

\textbf{Argument} --- a set of initial statements, called premises, followed by a conclusion.

\medskip

An argument is valid if and only if, in every case in which the premises are true, the conclusion is true. Otherwise, the argument is invalid.

\medskip

An argument can be valid if the conclusion is false. This can occur when a premise is false.

\bigskip

\textbf{Modus Ponens}

\medskip

If a conditional statement is true and the hypothesis of the conditional is true, the conclusion of the conditional statement is true.

\medskip

\begin{center}
\begin{tabular}{cc|c|c}
$p$ & $q$ & $p \to q$ & $\bigl(p \land (p \to q)\bigr) \to q$ \\
\hline
F & F & T & T \\
F & T & T & T \\
T & F & F & T \\
T & T & T & T
\end{tabular}
\end{center}

\bigskip

\textbf{Rules of Inference}

\medskip

\begin{center}
\begin{tabular}{@{}p{0.44\textwidth}l@{}}
$\displaystyle \frac{\begin{gathered}p\\p \to q\end{gathered}}{\therefore\ q}$ & \textbf{Modus Ponens} \\[1.3em]
$\displaystyle \frac{\begin{gathered}\lnot q\\p \to q\end{gathered}}{\therefore\ \lnot p}$ & \textbf{Modus Tollens} \\[1.3em]
$\displaystyle \frac{p}{\therefore\ p \lor q}$ & \textbf{Addition} \\[1.3em]
$\displaystyle \frac{p \land q}{\therefore\ p}$ & \textbf{Simplification}
\end{tabular}
\end{center}

\newpage

\textbf{Rules of Inference (continued)}

\medskip

\begin{center}
\begin{tabular}{@{}p{0.44\textwidth}l@{}}
$\displaystyle \frac{\begin{gathered}p \qquad q\end{gathered}}{\therefore\ p \land q}$ & \textbf{Conjunction} \\[1.3em]
$\displaystyle \frac{\begin{gathered}p \to q\\q \to r\end{gathered}}{\therefore\ p \to r}$ & \textbf{Hypothetical Syllogism} \\[1.3em]
$\displaystyle \frac{\begin{gathered}p \lor q\\\lnot p\end{gathered}}{\therefore\ q}$ & \textbf{Disjunctive Syllogism} \\[1.3em]
$\displaystyle \frac{\begin{gathered}p \lor q\\\lnot p \lor r\end{gathered}}{\therefore\ q \lor r}$ & \textbf{Resolution}
\end{tabular}
\end{center}

\bigskip

\textbf{Fallacy.} A fallacy resembles a rule of inference, but is based on contingencies instead of tautologies.

\medskip

Example: $\bigl((p \to q) \land q\bigr) \to p$.

\medskip

\begin{center}
\begin{tabular}{cc|c}
$p$ & $q$ & $\bigl((p \to q) \land q\bigr) \to p$ \\
\hline
F & F & T \\
F & T & F \\
T & F & T \\
T & T & T
\end{tabular}
\end{center}

\newpage

\bigskip

\textbf{Rules of Inference for Quantified Statements}

\medskip

\textbf{Universal Instantiation}

\medskip

For example, from the statement ``All clerics are wise,'' we can conclude that ``Kristen is wise'' when Kristen is a member of the domain of all clerics.

\medskip

\textbf{Universal Generalization}

\medskip

To show that $\forall x\,P(x)$ is true, take an arbitrary element $c$ from the domain and show that $P(c)$ is true. The element $c$ must be arbitrary.

\medskip

\textbf{Existential Instantiation}

\medskip

Existential instantiation is the rule that allows us to conclude that there is an element $c$ in the domain for which $P(c)$ is true if we know that $\exists x\,P(x)$ is true.

\medskip

\textbf{Existential Generalization}

\medskip

Existential generalization is the rule of inference used to conclude that $\exists x\,P(x)$ is true when a particular element $c$ with $P(c)$ true is known.

\medskip

\begin{center}
\begin{tabular}{@{}p{0.52\textwidth}l@{}}
$\displaystyle \frac{\forall x\,P(x)}{\therefore\ P(c)}$ & \textbf{Universal Instantiation} \\[1.5em]
$\displaystyle \frac{P(c)\ \text{for an arbitrary }c}{\therefore\ \forall x\,P(x)}$ & \textbf{Universal Generalization} \\[1.5em]
$\displaystyle \frac{\exists x\,P(x)}{\therefore\ P(c)\ \text{for some element }c}$ & \textbf{Existential Instantiation} \\[1.5em]
$\displaystyle \frac{P(c)\ \text{for some element }c}{\therefore\ \exists x\,P(x)}$ & \textbf{Existential Generalization}
\end{tabular}
\end{center}

\end{document}
